\documentclass[11pt,a4paper]{article}
\usepackage[utf8]{inputenc}
\usepackage[T1]{fontenc}
\usepackage[spanish,es-noquoting,es-noshorthands]{babel}
\usepackage{amsmath}
\usepackage[margin=2.2cm]{geometry}
\usepackage{listings}
\usepackage{xcolor}
\usepackage{booktabs}
\usepackage{longtable}
\usepackage[hidelinks]{hyperref}
\setlength{\parskip}{0.55em}
\setlength{\parindent}{0pt}

\definecolor{cmt}{rgb}{0.35,0.50,0.35}
\definecolor{kw}{rgb}{0.10,0.25,0.60}
\definecolor{str}{rgb}{0.60,0.20,0.20}
\definecolor{bg}{rgb}{0.975,0.975,0.97}

\lstset{
  language=C++,
  basicstyle=\ttfamily\scriptsize,
  commentstyle=\color{cmt}\itshape,
  keywordstyle=\color{kw}\bfseries,
  stringstyle=\color{str},
  backgroundcolor=\color{bg},
  numbers=left, numberstyle=\tiny\color{gray}, numbersep=7pt,
  breaklines=true, breakatwhitespace=false,
  showstringspaces=false, tabsize=4,
  frame=leftline, framesep=6pt, rulecolor=\color{gray},
  columns=fullflexible, keepspaces=true,
  captionpos=b,
}


\title{\textbf{Recorrido del código}\\[4pt]\large\texttt{highOrderElectroActivationFoamImplicitPETScDistributed}}

\author{}

\date{Generado el 2026-08-14 a partir de \texttt{highOrderElectroActivationFoamImplicitPETScDistributed.C} (8145 líneas)}

\begin{document}

\maketitle

\noindent\textbf{Cómo leer esto.} Los listados NO son copias: salen de \texttt{\textbackslash lstinputlisting} apuntando al fuente real, con rangos que este reporte recalcula al regenerarse. El PDF de hoy muestra el código de hoy. Los números de línea del margen son los del archivo.

\tableofcontents

\newpage

\section{Qué resuelve este solver}

Resuelve la propagación del potencial de acción cardíaco: la ecuación de
monodominio acoplada al modelo iónico \textbf{ten Tusscher--Noble--Noble--Panfilov
(TNNP)}, con el objetivo de reproducir el benchmark de Niederer et al. (2012).

A diferencia del solver MMS, acá \textbf{no hay solución exacta}. La verificación
es por convergencia, por invariancia frente a la descomposición, y contra
referencia publicada.

En cada paso de tiempo hay dos trabajos muy distintos:
%
\begin{enumerate}
  \item integrar, en cada celda o punto de Gauss, un sistema \textbf{rígido} de
        19 EDO --- el modelo TNNP --- para obtener la corriente iónica $I_{ion}$;
  \item resolver el sistema lineal implícito de la difusión.
\end{enumerate}
%
El primero domina el tiempo: la resolución lineal completa es apenas el 17--44\,\%
del lazo. Eso explica varias decisiones del archivo, entre ellas que la única
región OpenMP esté sobre las EDO y no sobre el ensamblado.

\textbf{Dos ejes de paralelismo}, en lugares distintos del archivo: \textbf{MPI}
reparte la malla entre ranks y toca todas las fases; \textbf{OpenMP} reparte el
lazo de puntos de integración dentro de un rank y sólo toca las EDO.

\section{El modelo iónico TNNP}

El modelo vive en una clase dentro del archivo, con sus 19 estados y sus
variables algebraicas.

\texttt{computeRates} evalúa las derivadas; \texttt{computeVariables} las
cantidades algebraicas; \texttt{calculateCurrent} devuelve $I_{ion}$.

\texttt{protectState} es la protección numérica: acota compuertas y
concentraciones a rangos físicos. Lleva un contador de correcciones, y ése es el
único estado compartido de todo el camino de las EDO --- por eso su incremento
va bajo \texttt{omp critical} y no \texttt{atomic}: el test de cota y el
incremento tienen que ser un solo paso indivisible.

El test \textbf{E01} verifica este modelo aislado, en 0-D, contra los valores
publicados de ten Tusscher 2004 --- APD90, pico, reposo, dV/dt máximo. Es el
test más barato de la suite y sostiene todo lo demás: si la corriente iónica
estuviera mal, ningún test de convergencia lo detectaría, porque todos
convergerían limpiamente a la respuesta equivocada.

\lstinputlisting[firstnumber=1973, firstline=1973, lastline=2029, caption={La protección numérica y su contador compartido.}]{../highOrderElectroActivationFoamImplicitPETScDistributed.C}%[firstnumber=1973, firstline=1973, lastline=2029, caption={La protección numérica y su contador compartido.}]


\lstinputlisting[firstnumber=2512, firstline=2512, lastline=2550, caption={La corriente iónica. Recortada.}]{../highOrderElectroActivationFoamImplicitPETScDistributed.C}%[firstnumber=2512, firstline=2512, lastline=2550, caption={La corriente iónica. Recortada.}]


\section{Los integradores de EDO, y el único paralelismo por hilos}

El solver trae sus propios steppers en vez de usar el marco \texttt{ODESolver}
de OpenFOAM: \texttt{eulerStep}, \texttt{rk4Step} y \texttt{rkf45Trial}
adaptativo. Un nombre desconocido es \texttt{FatalError}, no un fallback
silencioso.

Euler y RK4 dan \textbf{un} paso del tamaño del paso exterior completo, así que
su exactitud la controla \texttt{dt}; RKF45 subdivide por su cuenta y la
controla su tolerancia. El test \textbf{E02} hace converger los tres al mismo
potencial de acción, que es una afirmación mucho más fuerte que apretar uno
solo contra sí mismo.

El paso adaptativo \textbf{persiste entre llamadas} en \texttt{stepMs\_}, de modo
que cada llamada hereda el $h$ de la anterior. Eso hacía que la clave
\texttt{initialODEStep} fuera inalcanzable --- el constructor siembra
\texttt{stepMs\_} positivo, así que la rama que la leía sólo podía dispararse con
\texttt{deltaT == 0}. El test \textbf{E03} lo midió (cero efecto sobre cuatro
órdenes de magnitud) y la clave se removió.

\texttt{solveODE} contiene la \textbf{única} región OpenMP del archivo. Los
hilos reparten los puntos de integración, y funciona porque cada punto resuelve
su propia EDO sin hablar con ningún otro --- por eso el resultado es
bit-idéntico con cualquier cantidad de hilos, cosa que \textbf{E20} mide dando
exactamente cero. \textbf{El ensamblado no está paralelizado con hilos.}

\lstinputlisting[firstnumber=2138, firstline=2138, lastline=2175, caption={El paso adaptativo. Recortado.}]{../highOrderElectroActivationFoamImplicitPETScDistributed.C}%[firstnumber=2138, firstline=2138, lastline=2175, caption={El paso adaptativo. Recortado.}]


\lstinputlisting[firstnumber=2552, firstline=2552, lastline=2611, caption={El despacho de stepper y la región OpenMP. Recortado. \textit{(recortado; sigue en el fuente.)}}]{../highOrderElectroActivationFoamImplicitPETScDistributed.C}%[firstnumber=2552, firstline=2552, lastline=2611, caption={El despacho de stepper y la región OpenMP. Recortado. \textit{(recortado; sigue en el fuente.)}}]


\section{El residuo de estados: una colectiva que colgaba}

\texttt{maxStateRelativeL2Difference} da el cambio relativo máximo de los
estados entre dos iteraciones no lineales, y es parte del criterio de parada.

Es una función \textbf{colectiva}: reduce dos veces por estado. Y ahí había un
cuelgue. La versión anterior devolvía temprano cuando el arreglo de estados de un
rank venía vacío --- exactamente lo que le pasa a un rank cuya partición no tiene
celdas --- así que ese rank salía sin participar de las reducciones mientras los
demás quedaban bloqueados esperando. \textbf{No es un número mal calculado: es
que la corrida no termina.}

Está corregido de dos formas, porque cualquiera de las dos sola dejaba el cuelgue
vivo: el rank vacío ahora cae hasta los lazos de acumulación (que iteran cero
veces) y llega igual a los dos \texttt{reduce}, y el conteo de vueltas del lazo
se reduce con \texttt{minOp}, porque salía de \texttt{current[0].size()}, que es
rank-local y ni siquiera se puede leer en un rank vacío.

El test \textbf{E19} lo mide sobre una malla de dos celdas: antes moría en el
timeout de 300 s a np=4 y np=8, ahora corre en 1 s, igual que el control serial.

\lstinputlisting[firstnumber=3097, firstline=3097, lastline=3156, caption={La colectiva, con las dos correcciones. \textit{(recortado; sigue en el fuente.)}}]{../highOrderElectroActivationFoamImplicitPETScDistributed.C}%[firstnumber=3097, firstline=3097, lastline=3156, caption={La colectiva, con las dos correcciones. \textit{(recortado; sigue en el fuente.)}}]


\section{Bloques de construcción del ensamblado}

Antes de las matrices de rigidez vienen las piezas compartidas.

\texttt{addTripletIfNeeded} es el punto por el que pasa toda entrada de la
matriz, y descarta las menores que \texttt{SMALL}.

\texttt{addCellGradientDotCoeffs} agrega la fila de gradiente de una celda
proyectada sobre un vector, usando los stencils MLS. Es donde vive el problema
abierto más grande: sobre celdas pegadas a un corte de procesador ese stencil
pierde un miembro, y como el ajuste se rehace sobre el conjunto reducido,
\textbf{refita toda la fila}. Vive en solids4foam, no acá.

\texttt{exchangeCoupledFaceRows} intercambia con el rank vecino las filas que el
término de estabilización necesita del otro lado del corte. Es \textbf{colectiva}
y está deliberadamente \emph{fuera} del lazo de chunks: adentro se trabaría en
cuanto dos ranks tuvieran distinta cantidad de chunks.

\lstinputlisting[firstnumber=3457, firstline=3457, lastline=3469, caption={Por acá pasa toda entrada.}]{../highOrderElectroActivationFoamImplicitPETScDistributed.C}%[firstnumber=3457, firstline=3457, lastline=3469, caption={Por acá pasa toda entrada.}]


\lstinputlisting[firstnumber=3986, firstline=3986, lastline=4045, caption={El intercambio a través del corte. Recortado. \textit{(recortado; sigue en el fuente.)}}]{../highOrderElectroActivationFoamImplicitPETScDistributed.C}%[firstnumber=3986, firstline=3986, lastline=4045, caption={El intercambio a través del corte. Recortado. \textit{(recortado; sigue en el fuente.)}}]


\section{Las dos matrices de rigidez, y el defecto de la conductividad}

Igual que en el MMS hay dos ensamblados y la tríada elige cuál corre: el
ortogonal de dos puntos y el de alto orden sobre el stencil LRE.

Acá estuvo el defecto más grande que encontró la verificación. El camino de alto
orden usaba \texttt{conductivity[own]} \textbf{en las dos filas} de una cara
interna, en vez del promedio $\tfrac12(D_{own}+D_{nei})$ que ya usaban el camino
de bajo orden y el término de estabilización.

Con $D$ uniforme los dos tensores coinciden y la diferencia se cancela
idénticamente, así que fue invisible durante toda la vida del código: todas las
corridas usaban el tensor constante del benchmark. Con $D$ \textbf{no uniforme}
--- o sea, en cuanto entren fibras --- hacía que $K$ cambiara un \textbf{21\,\%}
al descomponer la malla, porque al cortar, cada rank pasa a ser dueño de su lado
de la cara y ninguno coincide con lo que hizo la corrida serial.

El arreglo son tres cambios, no uno: el promedio en la cara interna, el promedio
\textbf{también en la cara de procesador} (con el tensor remoto), y sacar el
intercambio de conductividad de la condición de estabilización, porque ahora hace
falta con cualquier alpha. Corregir sólo el primero arreglaba el operador serial
y dejaba el paralelo roto. \textbf{M15} lo mide: de 2.108e-01 a 4.736e-15.

\lstinputlisting[firstnumber=4559, firstline=4559, lastline=4618, caption={El ensamblado de alto orden. Recortado. \textit{(recortado; sigue en el fuente.)}}]{../highOrderElectroActivationFoamImplicitPETScDistributed.C}%[firstnumber=4559, firstline=4559, lastline=4618, caption={El ensamblado de alto orden. Recortado. \textit{(recortado; sigue en el fuente.)}}]


\section{Modo híbrido: EDO donde el frente lo pide}

Resolver la EDO en cada punto de Gauss es lo más directo y lo más caro;
resolverla en el centro de celda y reconstruir a los puntos es mucho más barato
pero falla donde el frente es abrupto, porque los estados del TNNP son rígidos.

\texttt{flagFrontCells} clasifica las celdas del frente por el rango de $V_m$
entre sus puntos de Gauss, y \texttt{dilateFrontCells} agranda esa banda unas
capas. El modo híbrido resuelve la EDO en los puntos de Gauss ahí y en el centro
de celda en el resto.

El detalle de implementación que importa: la clasificación se calcula
\textbf{una vez por paso} y se congela para todo el solve no lineal. Recalcularla
por iteración hace que el interruptor frente/suave oscile entre iteraciones y
convierte el punto fijo de Picard en un ciclo límite que nunca converge.

\lstinputlisting[firstnumber=5603, firstline=5603, lastline=5629, caption={La clasificación del frente. Recortada.}]{../highOrderElectroActivationFoamImplicitPETScDistributed.C}%[firstnumber=5603, firstline=5603, lastline=5629, caption={La clasificación del frente. Recortada.}]


\lstinputlisting[firstnumber=5640, firstline=5640, lastline=5699, caption={La dilatación de la banda. Recortada. \textit{(recortado; sigue en el fuente.)}}]{../highOrderElectroActivationFoamImplicitPETScDistributed.C}%[firstnumber=5640, firstline=5640, lastline=5699, caption={La dilatación de la banda. Recortada. \textit{(recortado; sigue en el fuente.)}}]


\section{Tiempos de activación y muestreo del benchmark}

Lo que el benchmark reporta son \textbf{tiempos de activación}, no campos.

\texttt{updateActivationTimes} los interpola \emph{dentro} del paso ---
$t_n + w\,\Delta t$ con $w$ el cruce lineal del umbral --- así que no quedan
cuantizados a múltiplos de $\Delta t$, lo cual importa para cualquier estudio de
refinamiento.

\texttt{writeActivationSamples} escribe el valor en el punto P8 y un perfil a lo
largo de la diagonal. El muestreo es en puntos \textbf{físicos fijos},
independientes de la malla, y eso es lo que permite restar dos perfiles índice a
índice sin interpolar entre discretizaciones --- la base de E11 y de los estudios
de convergencia.

\texttt{nearestCellToPoint} merece atención: devuelve el índice en \textbf{exactamente
un} rank y $-1$ en todos los demás, resuelto con una reducción para que la
elección sea determinista. El consumidor está guardado contra ese $-1$; el
\emph{banner} que imprime el punto de parada no lo estaba, y evaluaba
\texttt{mesh.C()[-1]} en todos los ranks perdedores. En build optimizado eso es
lectura fuera de rango, o sea comportamiento indefinido, y se manifestaba como
\textbf{SIGFPE} de forma reproducible pero aparentemente arbitraria: a dx=0.125
mm fallaba con 4 ranks y no con 1, 2 ni 8.

\lstinputlisting[firstnumber=5766, firstline=5766, lastline=5803, caption={La interpolación dentro del paso. Recortada.}]{../highOrderElectroActivationFoamImplicitPETScDistributed.C}%[firstnumber=5766, firstline=5766, lastline=5803, caption={La interpolación dentro del paso. Recortada.}]


\lstinputlisting[firstnumber=5807, firstline=5807, lastline=5866, caption={Exactamente un rank gana; el resto devuelve -1. \textit{(recortado; sigue en el fuente.)}}]{../highOrderElectroActivationFoamImplicitPETScDistributed.C}%[firstnumber=5807, firstline=5807, lastline=5866, caption={Exactamente un rank gana; el resto devuelve -1. \textit{(recortado; sigue en el fuente.)}}]


\section{main(): montaje y lazo temporal}

\texttt{main()} lee malla y diccionarios, construye los operadores LRE, ensambla
$M$ y $K$ \textbf{una sola vez} --- la no linealidad está en el término iónico,
no en la difusión --- y entra al lazo temporal.

Cada paso resuelve un sistema no lineal con Picard, JFNK o diagonalIion, y
\textbf{E10} verifica que los tres lleguen al mismo lugar. Un nombre de método
inválido es \texttt{FatalError} y no un fallback silencioso a Picard, cosa que
\textbf{E18} comprueba.

Al agotar las iteraciones sin converger, la rama está \textbf{una sola vez},
después del despacho de método, así que los tres comparten la semántica: por
defecto se descarta el paso y se vuelve a los valores del inicio, y con
\texttt{nonlinearAcceptUnconverged} se acepta el iterado. \textbf{E24} lo
verifica en las dos configuraciones. El solver MMS tiene el corte adentro del
despacho y por eso allá Picard acepta mientras los otros dos descartan --- dos
semánticas en la misma matriz de métodos, y sigue abierto.

\lstinputlisting[firstnumber=2767, firstline=2767, lastline=2791, caption={El paso estable de referencia.}]{../highOrderElectroActivationFoamImplicitPETScDistributed.C}%[firstnumber=2767, firstline=2767, lastline=2791, caption={El paso estable de referencia.}]


\lstinputlisting[firstnumber=2823, firstline=2823, lastline=2852, caption={El estímulo a nivel tejido. Recortado.}]{../highOrderElectroActivationFoamImplicitPETScDistributed.C}%[firstnumber=2823, firstline=2823, lastline=2852, caption={El estímulo a nivel tejido. Recortado.}]


\section{Índice de funciones}
Todas las funciones del espacio de nombres anónimo, en orden de aparición. Las secciones anteriores recorren las principales; ésta es el mapa completo para ubicarse en el archivo.
\begin{longtable}{rll}
\toprule
línea & función & líneas \\
\midrule
\endhead
94 & \texttt{gCol} & 4 \\
100 & \texttt{gNCols} & 4 \\
148 & \texttt{checkPetscError} & 10 \\
170 & \texttt{buildDistributedMat} & 85 \\
287 & \texttt{clear} & 15 \\
303 & \texttt{reset} & 19 \\
370 & \texttt{gSquaredNorm} & 6 \\
379 & \texttt{gNorm} & 4 \\
387 & \texttt{gDot} & 6 \\
396 & \texttt{lowerWord} & 12 \\
411 & \texttt{usesPetscBackend} & 5 \\
422 & \texttt{petscKspTypeName} & 25 \\
452 & \texttt{petscPcTypeName} & 43 \\
506 & \texttt{petscParallelPcTypeName} & 13 \\
522 & \texttt{isGmresType} & 4 \\
581 & \texttt{copyEigVecToPetscVec} & 12 \\
595 & \texttt{copyPetscVecToEigVec} & 16 \\
615 & \texttt{setPetscOptionsPrefix} & 11 \\
808 & \texttt{updateValues} & 62 \\
871 & \texttt{solve} & 42 \\
926 & \texttt{petscShellMatMult} & 17 \\
955 & \texttt{petscShellPCApply} & 17 \\
1037 & \texttt{initialise} & 112 \\
1499 & \texttt{embeddedComputeStimulus} & 43 \\
1543 & \texttt{computeIstim} & 14 \\
1911 & \texttt{tissueFlag} & 14 \\
1926 & \texttt{stateIndex} & 8 \\
1935 & \texttt{algebraicIndex} & 8 \\
1944 & \texttt{loadStimulusProtocol} & 10 \\
1960 & \texttt{recordCorrection} & 12 \\
1973 & \texttt{protectState} & 57 \\
2031 & \texttt{computeRates} & 40 \\
2072 & \texttt{computeVariables} & 19 \\
2092 & \texttt{eulerStep} & 16 \\
2109 & \texttt{rk4Step} & 28 \\
2138 & \texttt{rkf45Trial} & 38 \\
2177 & \texttt{advancePoint} & 133 \\
2316 & \texttt{emptyScratchIDs} & 5 \\
2469 & \texttt{initialiseStates} & 32 \\
2502 & \texttt{updateStatesOld} & 4 \\
2507 & \texttt{resetStatesToStatesOld} & 4 \\
2512 & \texttt{calculateCurrent} & 39 \\
2552 & \texttt{solveODE} & 108 \\
2661 & \texttt{exportStates} & 31 \\
2680 & \texttt{if} & 4 \\
2693 & \texttt{exportStatesIntegrationPoints} & 38 \\
2738 & \texttt{characteristicDx} & 24 \\
2767 & \texttt{computeStableDeltaT} & 25 \\
2795 & \texttt{currentPeakRSSMB} & 14 \\
2813 & \texttt{logMemoryCheckpoint} & 5 \\
2823 & \texttt{applyStimulus} & 30 \\
2863 & \texttt{computeHighOrderLaplacian} & 80 \\
2946 & \texttt{quadraticForm} & 10 \\
2964 & \texttt{reconstructFromTaylor} & 24 \\
2992 & \texttt{thetaFromScheme} & 18 \\
3013 & \texttt{normalizedMassMatrixType} & 18 \\
3035 & \texttt{fieldToEigVec} & 12 \\
3051 & \texttt{eigVecToField} & 9 \\
3063 & \texttt{relativeL2Difference} & 21 \\
3087 & \texttt{relativeL2Norm} & 5 \\
3097 & \texttt{maxStateRelativeL2Difference} & 66 \\
3176 & \texttt{diagonalBlock} & 23 \\
3203 & \texttt{addDiagonalToMatrix} & 18 \\
3235 & \texttt{solveGMRES} & 142 \\
3390 & \texttt{postProcessingDir} & 11 \\
3416 & \texttt{reportOperatorFingerprint} & 35 \\
3457 & \texttt{addTripletIfNeeded} & 13 \\
3473 & \texttt{assembleDiagonalMassMatrix} & 31 \\
3517 & \texttt{assembleConsistentMassMatrixHO} & 141 \\
3663 & \texttt{orthogonalDiffusionCoeff} & 14 \\
3680 & \texttt{normalDiffusivity} & 4 \\
3690 & \texttt{stabilisationFaceCoeff} & 17 \\
3711 & \texttt{addCellGradientDotCoeffs} & 58 \\
3786 & \texttt{cellTaylorEntryCoeff} & 32 \\
3827 & \texttt{addCellTaylorExtrapolationCoeffs} & 57 \\
3888 & \texttt{buildCellTaylorExtrapolationRow} & 29 \\
3921 & \texttt{buildCellGradientDotRow} & 26 \\
3951 & \texttt{addRemoteRowCoeffs} & 19 \\
3986 & \texttt{exchangeCoupledFaceRows} & 78 \\
4077 & \texttt{addRhieChowInternalJumpCoeffs} & 34 \\
4118 & \texttt{addRhieChowCoupledJumpCoeffs} & 38 \\
4161 & \texttt{addRhieChowBoundaryJumpCoeffs} & 21 \\
4194 & \texttt{assembleStandardOrthogonalStiffnessMatrix} & 345 \\
4559 & \texttt{assembleHighOrderStiffnessMatrix} & 478 \\
5046 & \texttt{solveSparseSystemEigen} & 94 \\
5149 & \texttt{solveSparseSystem} & 91 \\
5250 & \texttt{reconstructVmAtIionIntegrationPoints} & 80 \\
5340 & \texttt{reconstructStatesAtIionIntegrationPoints} & 138 \\
5492 & \texttt{clampPhysical} & 13 \\
5509 & \texttt{averageIionIntegrationPointsToCells} & 34 \\
5548 & \texttt{averageStatesIntegrationPointsToCells} & 47 \\
5603 & \texttt{flagFrontCells} & 27 \\
5640 & \texttt{dilateFrontCells} & 116 \\
5766 & \texttt{updateActivationTimes} & 38 \\
5807 & \texttt{nearestCellToPoint} & 62 \\
5875 & \texttt{sampleIDW} & 146 \\
6025 & \texttt{writeActivationSamples} & 63 \\
6090 & \texttt{main} & 2055 \\
\bottomrule
\end{longtable}

\end{document}